
\section{纤维积}
\label{sec:a1-2}

下面定义仓西结构的纤维积\liftIndex{I317}。设 $Y_i$ 是光滑流形，$X_i$ 是带仓西结构
的空间，并沿用引理~\ref{lem:a1-11} 的记号
\[
 \bigl(V_{p_i}^i,E_{p_i}^i,\Gamma_{p_i}^i,\psi_{p_i}^i,s_{p_i}^i;
 V_{p_iq_i}^i,\widehat\phi_{p_iq_i}^i,
 \phi_{p_iq_i}^i,h_{p_iq_i}^i\bigr),
 \qquad i=1,\ldots,I.
\]
令 $\{f_{i,p_i}\}\colon X_i\to Y_i$ 是光滑、强连续映射，并假设它们
弱浸没。置 $Y=\prod_iY_i$。设 $W$ 是带角流形，
$f\colon W\to Y$ 是光滑映射；又令 $f_i\colon X_i\to Y_i$ 是由
$\{f_{i,p_i}\}$ 诱导的连续映射。置
\begin{equation}
\begin{aligned}
 Z
 &=\left(\prod_iX_i\right)\times_YW
 \\
 &:=\left\{((p_i),p)\in\left(\prod_iX_i\right)\times W\ \middle|\
 \bigl(f_1(p_1),\ldots,f_I(p_I)\bigr)=f(p)\right\}.
\end{aligned}
\tag{A1.36}\label{eq:a1-36}
\end{equation}
下面在 $Z$ 上定义仓西结构。

\begin{definition}{A1.37}\label{def:a1-37}
设 $\vec p=((p_i),p)\in Z$，并令 $V_{p_i}^i$ 是 $p_i$ 的仓西邻域。
置
\begin{equation}
\begin{aligned}
 V_{\vec p}:={}&\left\{((x_i),x)\in
 \left(\prod_iV_{p_i}^i\right)\times W\ \middle|\right.\\[-0.2em]
 &\hspace{7.2em}\left.
 \bigl(f_{1,p_1}(x_1),\ldots,f_{I,p_I}(x_I)\bigr)=f(x)
 \right\}.
\end{aligned}
\tag{A1.38}\label{eq:a1-38}
\end{equation}
由于 $f_{i,p_i}$ 都是浸没，每个 $V_{\vec p}$ 都是光滑流形\translatornote{式 (A1.36)、(A1.38) 的源书末项分别写作
$f_n(p_n)$、$f_{n,p_n}(x_n)$；但本段只定义了指标范围
$i=1,\ldots,I$，故依同一元组的其余项恢复为第 $I$ 项。}。置
\[
 E_{\vec p}=\prod_iE_{p_i}^i,
 \qquad
 \Gamma_{\vec p}=\prod_i\Gamma_{p_i}^i.
\]
容易看出，$\Gamma_{\vec p}$ 作用在
$E_{\vec p}\times V_{\vec p}\to V_{\vec p}$ 上。利用各
$s_{p_i}^i$，可以按显然的方式定义映射
$s_{\vec p}\colon V_{\vec p}\to E_{\vec p}$；各 $\psi_{p_i}^i$ 则诱导
映射
\[
 \psi_{\vec p}\colon s_{\vec p}^{-1}(0)/\Gamma_{\vec p}\longrightarrow Z.
\]
直接验证可知，
$(V_{\vec p},E_{\vec p},\Gamma_{\vec p},\psi_{\vec p},s_{\vec p})$
是 $\vec p$ 的仓西邻域，并且这些邻域胶合成 $Z$ 上的仓西结构。
\end{definition}

\begin{lemma}{A1.39}\label{lem:a1-39}
若 $X_i$ 上的仓西结构都有切丛，则 $Z$ 上的仓西结构也有切丛。若
$X_i$ 上的仓西结构以及流形 $Y_i,W$ 都已定向，则 $Z$ 上的仓西结构
也已定向。
\end{lemma}

\begin{proof}
设 $q_i<p_i$，并取
\[
 x_i\in
 \psi_{p_i}^i\bigl((s_{p_i}^i)^{-1}(0)/\Gamma_{p_i}^i\bigr)
 \cap
 \psi_{q_i}^i\bigl((s_{q_i}^i)^{-1}(0)/\Gamma_{q_i}^i\bigr)
 \subset X_i.
\]
写成
\[
 x_i=\psi_{p_i}^i\bigl(\pi(x_{i,p_i})\bigr)
     =\psi_{q_i}^i\bigl(\pi(x_{i,q_i})\bigr).
\]
假设
\[
 \vec q=((q_i),q),
 \qquad \vec x=((x_i),x),
 \qquad \vec p=((p_i),p)
\]
都是 $Z$ 的元素。由式~\eqref{eq:a1-15}，仓西映射诱导同构
\begin{equation}
 \frac{T_{x_{i,p_i}}V_{p_i}^i}
 {(d_{x_{i,q_i}}\phi_{p_iq_i}^i)(T_{x_{i,q_i}}V_{q_i}^i)}
 \cong
 \frac{E_{p_i}^i}
 {\widehat\phi_{p_iq_i,x_i}^i
  (E_{q_i}^i|_{V_{p_iq_i}^i})}.
\tag{A1.40}\label{eq:a1-40}
\end{equation}
置
\[
 \vec x_{\vec p}=((x_{i,p_i}),x)\in V_{\vec p},
 \qquad
 \vec x_{\vec q}=((x_{i,q_i}),x)\in V_{\vec q}.
\]
由于存在正合列
\[
 0\longrightarrow T_{\vec x_{\vec p}}V_{\vec p}
 \longrightarrow
 \left(\bigoplus_iT_{x_{i,p_i}}V_{p_i}^i\right)\oplus T_xW
 \longrightarrow\bigoplus_iT_{x_i}Y_i\longrightarrow0
\]
和
\[
 0\longrightarrow T_{\vec x_{\vec q}}V_{\vec q}
 \longrightarrow
 \left(\bigoplus_iT_{x_{i,q_i}}V_{q_i}^i\right)\oplus T_xW
 \longrightarrow\bigoplus_iT_{x_i}Y_i\longrightarrow0,
\]
故有同构
\begin{equation}
 \frac{T_{\vec x_{\vec p}}V_{\vec p}}
 {(d_{\vec x_{\vec q}}\phi_{\vec p\vec q})
   (T_{\vec x_{\vec q}}V_{\vec q})}
 \cong
 \bigoplus_i
 \frac{T_{x_{i,p_i}}V_{p_i}^i}
 {(d_{x_{i,q_i}}\phi_{p_iq_i}^i)(T_{x_{i,q_i}}V_{q_i}^i)}.
\tag{A1.41}\label{eq:a1-41}
\end{equation}
因而式~\eqref{eq:a1-40} 蕴含
\begin{equation}
\begin{aligned}
 \frac{T_{\vec x_{\vec p}}V_{\vec p}}
 {(d_{\vec x_{\vec q}}\phi_{\vec p\vec q})
   (T_{\vec x_{\vec q}}V_{\vec q})}
 &\cong
 \bigoplus_i
 \frac{E_{p_i}^i}
 {\widehat\phi_{p_iq_i,x_i}^i
  (E_{q_i}^i|_{V_{p_iq_i}^i})}
 \\
 &=
 \frac{E_{\vec p}}
 {\widehat\phi_{\vec p\vec q,\vec x}
  (E_{\vec q}|_{V_{\vec p\vec q}})}.
\end{aligned}
\tag{A1.42}\label{eq:a1-42}
\end{equation}
这一同构由仓西映射的微分诱导。由此证明了切丛的存在性。可定向性的
证明类似。
\end{proof}

注意，映射 $Z\to W$ 由一个强连续映射诱导；我们把它写成
\[
 \widehat f=(f_i)\times_YW.
\]

\begin{lemma}{A1.43}\label{lem:a1-43}
假设 $W$ 与 $Y$ 均为紧致、无边界且已定向\translatornote{源书把假设写成“$W$ 与 $Y$ 紧致、无边界且
orientation”，语法上漏去“已赋予定向”的谓语；译文依本式中的基本类、
庞加莱对偶及引理 A1.39 的定向假设补明。}的流形，并且
$\partial X_i=\varnothing$。则
\[
 f^*\bigl(
 PD(f_{1*}([X_1]))\cup\cdots\cup PD(f_{I*}([X_I]))
 \bigr)
 =\mathbin{\pm}PD(\widehat f_*([Z])).
\]
这里 $PD$ 表示庞加莱对偶。
\end{lemma}

证明很容易，留给读者。（本书不使用引理~\ref{lem:a1-43}。）

设 $f_i\colon X_i\to Y$（$i=1,2$）是强连续映射，并令
\[
 Z=X_1\times_YX_2
 =\{(x_1,x_2)\in X_1\times X_2\mid f_1(x_1)=f_2(x_2)\}.
\]
若 $X_1,X_2$ 有仓西结构，且 $f_1,f_2$ 弱浸没\liftIndexRangeStart{I382}，则可把 $Z$ 看成
\[
 (X_1\times X_2)_{f_1\times f_2}\!\times_{Y^2,\Delta}Y,
\]
其中 $Y\to Y^2$ 是 $y\mapsto(y,y)$。于是 $Z$ 有仓西结构。事实上，
只需假设 $f_1,f_2$ 中有一个弱浸没即可。

\begin{remark}{A1.44}\label{rem:a1-44}
\textup{(1)} 我们只在 $f_i\colon X_i\to Y$（$i=1,2$）都弱浸没时
定义纤维积。一般情形可以通过修改仓西结构约化到这一情形。具体地说，
可如下同时增大 $V_{p_j^{(i)}}^i$ 与 $E_{p_j^{(i)}}^i$，使
$f_{i,p_j^{(i)}}$ 成为浸没。

令
\[
 \left\{
 (V_{p_j^{(i)}}^i,E_{p_j^{(i)}}^i,\Gamma_{p_j^{(i)}}^i,
  \psi_{p_j^{(i)}}^i,s_{p_j^{(i)}}^i)
 \right\}_{p_j^{(i)}\in X_i},
 \qquad i=1,2,
\]
是两套仓西邻域系。考虑图像
\[
 \Gamma_{f_{i,p_j^{(i)}}}
 \subset V_{p_j^{(i)}}^i\times Y
\]
并取其管状邻域 $U(\Gamma_{f_{i,p_j^{(i)}}})$。图像的法丛
$N(\Gamma_{f_{i,p_j^{(i)}}})$ 之单位圆盘丛可与该管状邻域识别；它自然
同构于 $\operatorname{pr}_2^*TY$ 的相应限制。这里
$\operatorname{pr}_\ell$ 表示从 $V_{p_j^{(i)}}^i\times Y$ 到第 $\ell$
个因子的投影。于是
\[
 \left(
 U(\Gamma_{f_{i,p_j^{(i)}}}),
 \left.\operatorname{pr}_2^*TY\right|_{U(\Gamma_{f_{i,p_j^{(i)}}})},
 \Gamma_{p_j^{(i)}}^i,
 \psi_{p_j^{(i)}}^i\circ\operatorname{pr}_1,
 s_{p_j^{(i)}}^i\circ\operatorname{pr}_1\oplus s_{\mathrm{can}}
 \right)
\]
是 $p_j^{(i)}\in X_i$ 的一个仓西邻域。这里
$\Gamma_{p_j^{(i)}}^i$ 在第二个因子 $Y$ 上平凡作用，而
$s_{\mathrm{can}}$ 是经上述识别拉回到
$U(\Gamma_{f_{i,p_j^{(i)}}})$ 上的重言截面。由于
\[
 f_{i,p_j^{(i)}}
 =\operatorname{pr}_2\circ(\operatorname{id},f_{i,p_j^{(i)}}),
\]
且 $\operatorname{pr}_2$ 在 $U(\Gamma_{f_{i,p_j^{(i)}}})$ 上的限制显然
是浸没，所以可以使 $f_i$ 弱浸没。

也可以采用上述观察：纤维积 $X_1\times_YX_2$ 可识别为
\[
 Z=(X_1\times X_2)_{f_1\times f_2}\!\times_{Y\times Y,\Delta}Y,
\]
其中 $\Delta\colon Y\to Y\times Y$ 是对角嵌入。前面还指出，只要
$f_1\times f_2$ 或 $\Delta$ 之一弱浸没，这个纤维积就带有仓西结构。
以 $U(Y)$ 表示 $\Delta(Y)$ 在 $Y\times Y$ 中的一个管状邻域，以
$N_Y(Y\times Y)$ 表示其法丛。则
\[
 (U(Y),N_Y(Y\times Y),s_{\mathrm{can}})
\]
给出 $Y$ 的一个仓西结构，其中 $s_{\mathrm{can}}$ 是重言截面。注意，
$U(Y)\hookrightarrow Y\times Y$ 是开嵌入，特别是浸没。由此可见，
$\Delta$ 弱浸没，故 $Z$ 带有仓西结构\translatornote{源书把图像法丛写成管状邻域自身的
$N(U(\Gamma_{f_{i,p_j^{(i)}}}))$，但其单位圆盘丛应当给出
$U(\Gamma_{f_{i,p_j^{(i)}}})$，故按图像 $\Gamma_{f_{i,p_j^{(i)}}}$ 的
法丛及其拉回理解；末段又把开嵌入写成 $U(Y)\to Y$，其维数不合，
并且这里所需的底空间是 $Y\times Y$，故恢复目标为 $Y\times Y$。}。

\textup{(2)} 在适当假设下，把上述构造推广到 $Y$ 未必是流形而只是
带仓西结构空间的情形，看来是合理的。这样将得到一个范畴：其对象是带
仓西结构空间，其态射是某种强连续光滑映射。这个范畴会是概形或叠的范畴
的一种光滑类比。（我们指出，概形的纤维积总是良定义的，并且在概形理论
中起重要作用。）由于本书不用这样的范畴，我们不打算严格定义它。
\end{remark}
\liftIndexRangeEnd{I382}
